% Generated draft: requires mathematical and source review.
\documentclass[11pt]{article}
\usepackage[margin=25mm]{geometry}
\usepackage{fontspec}
\setmainfont{Latin Modern Roman}
\usepackage{amsmath,amssymb,mathtools,unicode-math}
\setmathfont{Latin Modern Math}
\usepackage{enumitem,needspace,xurl,hyperref}
\setlength{\parindent}{0pt}
\setlength{\parskip}{5pt}
\setlength{\emergencystretch}{3em}
\newcommand{\N}{\mathbb{N}}
\newcommand{\Z}{\mathbb{Z}}
\newcommand{\Q}{\mathbb{Q}}
\newcommand{\R}{\mathbb{R}}
\newcommand{\C}{\mathbb{C}}
\newcommand{\F}{\mathbb{F}}
\newcommand{\A}{\mathbb{A}}
\newcommand{\PP}{\mathbb{P}}
\newcommand{\E}{\mathbb{E}}
\newcommand{\eps}{\varepsilon}
\newcommand{\dd}{\,\mathrm{d}}
\begin{document}
\section*{UnsolvedMath ID: 10000025}
\textbf{Source problem code:} catalogue entry
\subsection{Definitions and mathematical statement}
For a graph \(G=(V,E)\) and finite \(A\subset V\), write \(\partial_E A=\{\{x,y\}\in E:x\in A,\ y\notin A\}\). One usual meaning of isoperimetric dimension greater than \(1\) is the existence of \(d>1\) and \(c>0\) such that
\[
 |\partial_E A|\ge c|A|^{(d-1)/d}\qquad(\varnothing\ne A\subset V\text{ finite}).
\]
For a vertex \(o\), let \(c_n(G,o)\) count length-\(n\) walks \((o=v_0,v_1,\ldots,v_n)\) with adjacent consecutive vertices and all vertices distinct. Define the upper rooted exponential growth rate by
\[
 \overline\mu(G,o)=\limsup_{n\to\infty}c_n(G,o)^{1/n}.
\]
When the limit exists, the rooted connective constant is \(\mu(G,o)=\lim_{n\to\infty}c_n(G,o)^{1/n}\). Conjecture 4.34 in the supplied record asks to prove that every graph \(G\) of isoperimetric dimension greater than \(1\) satisfies \(\mu(G)>1\). To make this exact, the source's graph class, boundary convention, and definition of \(\mu(G)\) must be supplied. In particular, the record does not say that \(G\) is connected, locally finite, of bounded degree, or transitive, does not specify a root, and does not ensure that a root-independent limit exists. Under the displayed edge-boundary interpretation and a rooted upper-growth interpretation, the corresponding candidate assertion is \(\overline\mu(G,o)>1\) for every vertex \(o\) of every admissible \(G\). This candidate is conditional on those conventions; it does not resolve the missing conventions of the source statement.
\subsection{Short English statement}
The conjecture links an isoperimetric inequality of dimension strictly greater than one to exponential growth of self-avoiding walks. The supplied statement leaves the graph class and connective-constant convention implicit. Its supplied open status is a triage assessment, with source verification dated 2026-07-29; no resolution is asserted here.
\subsection{Sources}
\begin{itemize}
\item Local UnsolvedMath registry, permanent ID \(10000025\), exact source code \texttt{catalogue entry}, ``Isoperimetric dimension and connective constants''; attributed to Itai Benjamini, no proposed year supplied.
\item Itai Benjamini, \emph{Coarse Geometry and Randomness}, Saint-Flour notes, Conjecture 4.34, PDF page 38: \url{https://www.wisdom.weizmann.ac.il/~itai/stflouraug24.pdf#page=38}. Supplied preservation URL: \url{https://arquivo.pt/noFrame/replay/20201231041548id_/http://www.wisdom.weizmann.ac.il/~itai/stflouraug24.pdf}.
\item Supplied provenance: accessed 2026-07-29; pypdf extraction checked against the rendered PDF; retained after cross-document deduplication and source-side status review. The source labels the item open or conjectural; the bounded modern audit found no complete primary-literature resolution. The summary mentions partial heuristics in the self-avoiding-walk community without a citation.
\item Dataset attribution: Ulam AI / UnsolvedMath Contributors, CC BY 4.0. The source-specific rights status is NEEDS\_REVIEW: publicly accessible author-hosted or publisher copy, with release reuse terms requiring review.
\end{itemize}
\end{document}
